\documentclass[a4paper]{article}
\usepackage[pdftex]{graphicx}
\usepackage[utf8]{inputenc}
\usepackage{enumerate}
\usepackage{siunitx}
\sisetup{locale=DE}
\usepackage{href-ul}
\hypersetup{
	colorlinks=true,
	linkcolor=blue,
	urlcolor=blue}
\usepackage{icomma}
\usepackage{amssymb}
\usepackage{geometry}
\geometry{a4paper, top=15mm, left=15mm, right=15mm, bottom=15mm,
	headsep=10mm, footskip=12mm}

\begin{document}
	%Title
	{\bf Exercises on strength of materials }\\
	\begin{enumerate}[1.]
		
		%Task 1
		\begin{minipage}{0.48\textwidth}
			\item To investigate the elasticity, test rods with the cross section $A=\SI{50}{\milli\meter^2}$ are loaded by attaching different masses and the resulting change in length is measured with a strain gauge. The following diagram is obtained for steel.
		\end{minipage}
		\hfill
		\begin{minipage}{0.48\textwidth}
			\includegraphics[width=5 cm]{emodul219}
		\end{minipage}
		
		\begin{enumerate}[a)]
			\item Determine using the given diagram
			\begin{enumerate}[(i)]
				\item the elastic modulus of steel
				\item is the mass at which the elastic deformation is $\varepsilon = \SI{10}{\micro\meter\over \meter}$.
				\item is the change in length in mm that a 2.5~m long steel rod with the same cross-section experiences at a tension of $\SI{4}{\newton\milli\meter^{-2}}$.
			\end{enumerate}
			\item In the given diagram, draw the graph for a material with an elastic modulus half as large
		\end{enumerate}
		
		%Exercise 2
		\begin{minipage}{0.48\textwidth}
			\item A steel pipe is installed without tension at $\SI{20}{\celsius}$ so that it cannot give way. It is heated to $\SI{80}{\celsius}$ by hot gases flowing through it.
		\end{minipage}
		\hspace{0.7 cm}
		\begin{minipage}{0.48\textwidth}
			\includegraphics[width=5 cm]{rohr219}
		\end{minipage}
		
		The linear expansion coefficient is $ \alpha=\SI{12e-6}{\kelvin^{-1}}$, the elastic modulus is $E=\SI{2.1e5}{\newton\milli\meter^{-2}} $
		\begin{enumerate}[a)]
			\item Calculate the compressive stress that occurs in the pipe.
			\item Calculate the pressure force in the pipe
		\end{enumerate}
		
		%Task 3
		\begin{minipage}{0.55\textwidth}
			\item A split flywheel is held together by two shrink rings that enclose both hub halves. The yield strength of the material is $\SI{285}{\newton\milli\meter^{-2}}$, the safety number should be $2.0$ and the two hub halves should be with the force $F=\SI{1100} {\kilo\newton}$ are held together.
		\end{minipage}
		\hfill
		\begin{minipage}{0.42\textwidth}
			\includegraphics[width=5 cm]{nabe219}
		\end{minipage}
		\begin{enumerate}[a)]
			\item Calculate the necessary width $b$ of the rings
			\item What inner diameter must the rings be made with?
		\end{enumerate}
		
		%Task 4
		\begin{minipage}{0.6\textwidth}
			\item A copper bushing and two rigid washers are secured by a central steel screw
			connected with each other. At the temperature $t_0=\SI{15}{\celsius}$ the system is installed without play but without preload. Calculate the stresses occurring in the bushing and screw when heated to $\SI{215}{\celsius}$. \\
			$[E_{Cu}=\SI{1,1e5}{\newton\milli\meter^{-2}}; \quad E_{St}=\SI{2,1e5}{\newton\milli\meter^{-2}}; \quad
			A_{St}=\SI{58}{\milli\meter^2}; \quad \alpha_{Cu}=\SI{1.65e-5}{\kelvin^{-1}}; \quad d_{ACu}=\SI{26}{\milli\meter};
			\quad d_{ICu}=\SI{20}{\milli\meter}]$
		\end{minipage}
		\hfill
		\begin{minipage}{0.38\textwidth}
			\includegraphics[width=4 cm]{stacu219}
		\end{minipage}
		
		%Task 5
		\item From a steel sheet with the tensile strength $\SI{340}{\newton\milli\meter^{-2}}<R_m<\SI{440}{\newton\milli\meter^{-2}}$ and With a thickness of 6.0~mm, discs with a diameter of 36~mm should be cut out. The shear stress is at most $80\%$ of the tensile strength. Calculate the maximum compressive stress acting in the punch of the cutting tool.
		
		%Task 6
		
	\end{enumerate}
	\fbox{
		\begin{minipage}{0.5\textwidth}
			Additional worksheets are available at
			
			\href{http://www.okuyakl.com}{www.okuyakl.com}
		\end{minipage}
		\hfill
		\begin{minipage}{0.4\textwidth}
			\includegraphics[width=1.5 cm]{../../viecher/zwe03}
			\includegraphics[width=2 cm]{../../viecher/afanticon1}
			
	\end{minipage}}
	%End
\end{document}